% ch2 B.4 后半 OPW 方法 (原书页 70–76 / PDF p085–p091)

% 衔接说明：本文件为小节 “O.P.W.（正交化平面波）方法” 的后半，不含小节标题。
% 首行承接原书 p.69（PDF p084）末句 “In general, the cancellation is less”，
% 本文件自 “complete for higher and higher l values.” 起续译（与 ch2_B4a.tex 末行拼接）。

% ---- 原书 p.70 (PDF p085) ----
完全——$l$ 值越高，就越是如此。对此有一个非常明显的原因：只要考察角动量为 $L$ 的轨道的径向波动方程
\begin{equation*}
\frac{\mathrm{d}^2 u_L}{\mathrm{d}r^2} - \left[\frac{L(L+1)}{r^2} + V(r)\right] u_L(r) = E_L u_L(r)
\end{equation*}
并勾画出其中的势能项（图 16）。我们看到，离心势项把束缚态从 $V$ 的内部区域排除出去，因此对较大的 $L$，束缚态无论如何也不可能为正是 $V$ 最强的那个区域本身构成一个完备集。尽管如此，虽然这个想法并没有明显地出现在数学表述之中，

%FIG{16}: 图 16　径向方程中势能项的草图：实线为 $V(r)$，虚线为离心势 $V + \frac{L(L+1)}{r^2}$，其谷底处标有 “Bound state?”（束缚态？）（原书图注仅作 “Figure 16”）

$V$ 的大矩阵元对这些较高的 $L$ 态实际上起不了多大作用，因为平面波态同样不能很有效地穿入芯区——离心项对它们同样起作用。这样看来，必然要发生的抵消之中，有一部分来自 $k$ 相当高的平面波，这也是这些态收敛性差的部分原因。事实上，Heine 曾指出，一定程度的抵消会显式地出现在久期方程之中。

我相信，这里的抵消理论仍然有进一步改进的余地。读者当中做核物理的人也许已经注意到，其中某些想法——无论是在 Q.D.M. 中还是在“被抵消的 O.P.W.”（cancelled O.P.W.）中——与多重散射理论（multiple scattering theory）和有效力程理论（effective range theory）的技术、以及在该理论中把散射矩阵 “T” 用作有效势的做法（38），有着惊人的相似。我想，通过与多重散射理论作类比，大概还能做成更多的事情，特别是把高动量的连续谱态连同束缚态一起从问题中消去。\footnote{在写完这些讲稿之后我获悉，Ham 与 Segall 的 Green 函数方法［Phys. Rev., \textbf{124}, 1786 (1961)］正是朝这个方向迈出的第一步：这一方法以及增广平面波方法（augmented plane-wave method）（M. M. Saffren，学位论文，M.I.T.，1959）之所以收敛极佳，很可能正是由于它们与多重散射理论的相似性。（译注：原文此处印作 “N.I.T.”，应为 M.I.T.。）}

% ---- 原书 p.71 (PDF p086) ----
用带抵消的 O.P.W. 所完成的详细定量工作是专家的课题。另一方面，我倒想提请各位注意它有助于说明的固体行为的若干定性特征。这些特征基本上有两类：（1）“好金属”行为中由它解释的那些方面；（2）固体量子化学中由 O.P.W. 何时良好、何时失效来解释的那些方面。

1. 这些方面中对剑桥的人士来说更熟悉的一个，是近自由电子图像在定性解释所观测到的纯金属费米面方面的显著成功。我想这方面已广为人知，无须进一步讨论。在这一推广的几个显著的失败之中，铜是一例：其中相关的芯能级——d 能级——下得不够深，无法用正交化来处理；而各价半导体（valence semiconductors）定量地看，离开自由电子球其实并不那么远。

第二个特征，是所谓“刚性能带”（rigid band）模型在无规合金和液态金属中的经验成功——正如 Heine 与 Ziman (39) 所强调的，这一成功实在可观。在这个模型中，人们忽略物质的无规背景结构，假装其行为受某种平均能带的控制，这个能带含有各组分所贡献的价电子的平均数目。在模型的极端形式下，人们假定金属的性质只由每原子的价电子数目决定，而不管这些电子由哪些原子贡献。在某些情形下，甚至晶体结构也能靠这种“对周期表取平均”的办法来预言。一个引人注目的例子是 Mo 与 Ru 的合金：它在晶体结构和其他若干性质上都模拟人造元素锝，而在所有这些性质上 Mo 和 Ru 都与 Tc 大相径庭 (40)。近来发现，越来越多的性质随电子/原子比这个参量相当平滑地变化，而这个参量当然只是决定相应费米面的大小 (10)。我不想造成一种错误印象，似乎刚性能带总能对合金与化合物的所有性质给出正确的描述——也许它们不奏效的时候比奏效的时候更多，特别是在混合彼此差别很大的金属时——而是想说：鉴于决定金属大部分物理性质的，是各种效应之间敏感的平衡，它们居然能够奏效，这件事本身就相当令人惊讶。

2. Phillips、Cohen 与 Heine 强调了抵消定理与实测出现的价晶体（valence crystals，有别于金属）之间的关系。在周期表中 $Z$ 较低的成员里，可以用来展开式 (18) 中 $\delta$ 函数的芯轨道少得多。这意味着，即便在较低的元素中，势的 $s$ 部分的抵消效率也较低；而在第一行中根本没有 p 抵消，d 抵消则要到第二长周期的末尾——从镓开始——才出现。于是，剩下的有效势 $(V + V_R)$ 将是周期表行数的减函数——例如从 C 到 Si 到 Ge 到 Sn 到 Pb 依次减小。因此，较早的成员将有更强的区界、更宽的禁带，以及形成开放价结构的更大的能量上的好处。所有这些预言都在第四列的这些元素中得到漂亮的例证。

至于随周期表第二个维度（即沿一行）的变化，则讨论得没有那么透彻。我想这是一个同样能被 O.P.W. 思想照亮的课题，我愿在此稍作讨论。

% ---- 原书 p.72 (PDF p087) ----
下面的讨论将基于这样一个事实：判定固体中价电子的波函数能否被相当好地描述为一组微扰效应相对较小的正交化平面波，实际上有两条判据。第一条我们已经见过——即芯能级必须相当深，使得正交化所给出的有效势的内芯部分已被相当好地抵消，而且芯能级不会强烈地混入价能级之中。

第二条判据可以这样来理解：从能量上看，固体中原子的电子组态要想与自由原子的相差悬殊，是完全不可能的。这一点当然已被实验充分证实：固体中原子的 X 射线结构因子与自由原子相比，从不曾有过挪动一两个以上电子那样严重的差异。各位当中也许有人注意过关于 X 射线测量的那场讨论：测量似乎表明铁在固体中少了四个 d 电子；虽然这些实验后来当然被证明是错的，但这件事至少促成了 Herring 的决定性论证 (41)：挪动这四个电子所需的能量，根本不可能来自内聚中所涉及的那几类能量。把两三个电子挪动一个与 Bohr 半径 $a_0$ 相比可观的数量——从一个壳层挪到另一个壳层——要付出若干 Rydberg 的能量。

因此，第二条重要判据是：\emph{我们预期价电子具有的那个原子波函数，必须能够相当好地用 O.P.W. 费米海内平面波态的占据来描述。}

碰巧，在某些条件下，原子态的构造确实不坏，足以按那种方式被近似。为说明这一点，我们可以去查原子波函数的 Fourier 变换，即它们在动量空间中的波函数［例如查 Bethe-Salpeter 的书，文献 (42)］。举几个例子：$1s$ 函数为
\begin{equation*}
\phi_{1s} = A\,\mathrm{e}^{-\kappa r}
\end{equation*}
而它的（未归一化的）Fourier 变换是
\begin{equation*}
\frac{\kappa}{\left(\kappa^2 + k^2\right)^2}
\end{equation*}
$2p$ 函数的则是
\begin{equation*}
\frac{\kappa^2 k\,Y_{1m}(k)}{\left(\kappa^2 + k^2\right)^3}
\end{equation*}
$3d$ 函数的则是

% ---- 原书 p.73 (PDF p088) ----
\begin{equation*}
\frac{\kappa^3 k^2}{\left(\kappa^2 + k^2\right)^4}\,Y_{2m}
\end{equation*}

这些函数的草图见图 17。当 $k$ 增大时，所有这些函数都在 $k = \kappa$ 附近相当陡峭地下降。如果费米面离这一点不太远，用 O.P.W. 便可对原子函数得到一个不错的近似。

%FIG{17}: 图 17　动量空间波函数 $\phi_k$ 随 $k$ 的变化：$1s$、$3d$、$2p$ 三条曲线均在 $k \approx \kappa$ 附近陡降（原书图注仅作 “Figure 17”）

如果考虑一个由原子函数 $\phi_{nl}$ 组成的 Bloch 波，便容易形象化地看出作出这种近似的第二条判据。对这样一个波，一个合理的近似可以是
\begin{equation*}
b_{\mathbf{k};nl} = \mathrm{e}^{i\mathbf{k}\cdot\mathbf{r}}\sum_j \phi_{nl}(\mathbf{r} - \mathbf{R}_j)
\end{equation*}
立即可见，这样一个函数的动量空间表示是
% 校注：原书求和号下印作小写 k；按物理意义（下页“与 k 由一个倒格矢相连的
% 每一个点”）求和指标应为倒格矢 K，此处按本意转写。
\begin{equation*}
b_{\mathbf{k};nl}(\mathbf{p}) = \sum_{\mathbf{K}} \delta_{\mathbf{p},\,\mathbf{k}+\mathbf{K}}\,\phi_{nl}(\mathbf{k} + \mathbf{K})
\end{equation*}

% ---- 原书 p.74 (PDF p089) ----
在与 $\mathbf{k}$ 由一个倒格矢相连的每一个点上，它取值 $\phi_{nl}(\mathbf{k} + \mathbf{K})$。以 $1s$ 函数为例在一维情形下把这一点形象化，我们可以把这个函数画成图 18 那样。在动量空间中，我们在 $\mathbf{k} + \mathbf{K}$ 的每个值处都有一个尖峰，其高度由原子动量空间波函数在该点处的大小给出。

%FIG{18}: 图 18　一维动量空间中的 $b(p)$：在由倒格矢与 $k$ 相连的各点 $k-K$、$0$、$k$、$k+K$ 处的尖峰，其高度由原子动量空间波函数 $\phi_{1s}$（钟形曲线）在该点的大小给出（原书图注仅作 “Figure 18”）

近自由电子模型中的波函数由单个平面波构成，充其量由少数几个能量十分接近的平面波构成。于是，比方说像上图中那样，如果 $\mathbf{k}$ 与 $\mathbf{k} - \mathbf{K}$ 都带有大的尖峰，而能量又不十分相近，这两类波函数就会非常不同。也就是说，$\mathbf{K}$ 矢量相对于 $k_F$ 不得太小，而 $k_F$ 必须 $\approx \kappa$。这等价于说：重要的布里渊区边界不得深入费米面之内太远，因为布里渊区边界恰恰就是 $|\mathbf{k} - \mathbf{K}|$ 等于 $|\mathbf{k}|$ 的地方，而远在 BZ 边界之外，$|\mathbf{k}|$ 便比 $|\mathbf{k} - \mathbf{K}|$ 大出相当多。毫无疑问，这一要求容许相当大的宽容度，但可以肯定，当
\begin{equation*}
|\mathbf{k} - \mathbf{K}| < \frac{3}{4}\,|\mathbf{k}|
\end{equation*}
左右时，这两类函数之间就不再有多少相似之处了。

这一要求的物理意义无非是：近自由电子图像无法很准确地描写变化剧烈的电子密度，因此，如果波函数的线度 $(1/\kappa)$ 比原子间距 $(1/K)$ 小，O.P.W. 就不可能是好的近似。换一个角度看：当 $1/\kappa$ 太小时，剩余的芯势开始变得太强。头一个要求在物理上同样不难理解：原子波函数中的动量分布与固体中的动量分布必须在大致相同的最大动能处截止；这就要求 $\kappa \approx k_F$。

% ---- 原书 p.75 (PDF p090) ----
这些要求中的头一条给了我们一个关于“好金属”行为的非常直截了当的判据；我们已经看到，波函数中的截止点与 $k_F$ 必须落在彼此非常靠近的地方，因而 BZ 不得深入费米球之内太远。但至少在 Bravais 的简单晶格中我们知道，第一布里渊区恰好含有每原子两个电子态，于是，若价电子数为 $Z$，则
\begin{equation*}
\text{BZ 的平均半径} = k_F \Big/ \left(\frac{Z}{2}\right)^{1/3}
\end{equation*}
（当 $Z = 2$ 时，BZ 的体积 = 费米球的体积。）因此这非常直接地就是对 $Z$ 本身的一个限制：它表明，$Z$ 大到 4 或 5 时还可以得到合理的拟合，但从此往后波函数将开始在径向方向上变化得非常快，平面波模型便行不通了。

看一看这些想法沿周期表的两个典型的行如何展开是很有教益的：一个是第二短周期，Na、Mg 等；另一个是第一过渡系，K、Ca、Si、Ti、……（译注：原文如此，此处 “Si” 应为 “Sc”。）

在 Na 等所在的一行中，我们先填充 $3s$ 能级，然后再填充 $3p$ 能级。$3s$ 与 $3p$ 原子能级在其动量空间表示中，正如在其径向空间函数中一样，带有径向节点（radial nodes）；但正交化过程将使恰好相应的节点结构出现在 O.P.W. 函数之中。因此，我们可以在一个“抹平”的空间中进行思考——其中径向节点被忽略——并假定波函数在这个“抹平”空间中的形状恰恰就是 $1s$ 或 $2p$ 函数的形状。

% 校注：原书此处指数印作 $r^{l'}$（l 后带一撇状墨点），按 Slater 波函数本意
% 转写为 $r^{\,l}$。
1930 年，Slater 从若干 Hartree 计算和各种经验数据出发加以推广，得出了一个非常成功的近似原子波函数与“屏蔽常数”（screening constants）(43) 的体系，供原子和分子计算使用。这些是抹平了的无节点函数，形式为 $r^{\,l}\,\mathrm{e}^{-\kappa r}$；而且看来，在我们的计算中所应使用的 $\kappa$ 多半恰好就是这些公式中的 $\kappa$。这些屏蔽常数和波函数，常被用作原子重要外区中的势以及相应波函数的一个相当粗糙的估计。我们在这里对这些元素给出：由金属中的原子体积算出的 $k_F$、由 Slater 按上述办法算出的 $\kappa$，以及 $\kappa_{1/2}$ 这个量——即原子动量函数降至其峰值约一半时的那个 $k$。见表 5。符合得并不完美（注意 Mg 和

\begin{table}[htbp]
\centering
{\bfseries 表~5}\\[5pt]
\begin{tabular}{lccccc}
\toprule
 & Na & Mg & Al & Si & P\\
\midrule
$\kappa$ & $1.39\times10^{8}$ & $1.80$ & $2.20$ & $2.60$ & $3.00$\\
$\kappa_{1/2}$ & $0.87$ & $1.12$ (p:1.6) & $2.00$ (s:1.4) & $2.32$ & \\
$k_F$ & $0.90$ & $1.36$ & $1.55$ & $1.84$ & \\
\bottomrule
\end{tabular}
\end{table}

% ---- 原书 p.76 (PDF p091) ----
\noindent Al 中的误差表明，在这两种情形下 s 与 p 的性质有所混合），但它一点也不坏；而且令人感到有趣的是：仅仅看一看原子波函数的形状，就能把固体的体积估计得多么漂亮。P 有各种各样复杂的晶体结构，一点也不像金属，也不遵守 $\kappa_{1/2} \cong k_F$ 的关系。因此，P、S 和 Cl 无法纳入 O.P.W. 理论：因为要使价电子分布相当均匀，所需的密度会高到使费米能从而平均动能远高于原子中的相应数值。换一种说法：在原子中，壳层简并（shell degeneracy）的存在使得一个填充得相当稠密的电子壳层可以在不付出过多动能的情况下被填满，因为不相容原理是凭借不同动量分量之间相当复杂的相位关系而得到满足的。这一效应只在壳层中电子数多于两三个的那些元素中才重要；但在这些元素中，它阻碍了真金属的形成，因为在真金属中这些相位关系必然要被破坏，而在芯附近维持给定电子密度所需的动能也就会太高。

在 P、S、Cl 中，正如在 N、O、F 中一样，其结果是形成共价键合的分子体系，其中 p 壳层的空穴在很大程度上仍保持其原子 p 函数的身份，并形成有方向的、饱和的共价键。

同样的过程也发生在过渡系中，只是结果不同。这里，所涉及的电子被正交化于 $(1s)^2$、$(2s)^2$、$(3s)^2$、$(2p)^6$ 和 $(3p)^6$。于是，单凭正交化本身，它们便可以带上 $4s$、$4p$ 和 $3d$ 的性质。$k$ 接近零的 O.P.W. 不可能带有可观的 d 性质，只能带有弱的 p 性质；因此，自然而然，起初在 K 中甚至在 Ca 中，O.P.W. 实际上就是 $4s$ 波函数。然而，在较高的 $k$ 值处，$3d$ 函数与 $4p$ 函数之间开始出现竞争。$4s$ 函数对 $4p$ 有有效的屏蔽，对 $3d$ 则没有；而较高的 $n$ 值意味着 $\kappa_{4p} < \kappa_{3d}$，于是在这一点上，$4p$ 函数在 $k$ 空间中的密度可以忽略，而电子——本质上仍是 O.P.W.——开始填充 $3d$ 能级。把屏蔽常数 $\kappa$ 与相应的费米动量加以比较即可表明这一点，见表 6。实际上，这里的符合甚至比前面还要好；而且显然，Sc、Ti 甚至 V 都是相当令人满意的 “O.P.W. 金属”。（括号中的数值是在如下假定下算出的：在这一系列的中段附近，
\begin{table}[htbp]
\centering
{\bfseries 表~6}\\[5pt]
\begin{tabular}{lccccccc}
\toprule
 & K & Ca & Sc & Ti & V & Cr & Mn\\
\midrule
$\kappa$ & $1.1\times10^{8}$ & $1.4$ & $1.65$ & $2.05$ & $2.45$ & $2.85$ & $3.25$ ($\kappa_{3d}$)\\
$\kappa_{1/2}$ & $0.68$ & $0.87$ & $1.65$ & $2.05$ & $2.45$ & $2.85$ & $3.25$ ($\kappa_{3d}$)\\
 & & & & & $(2.2)$ & $(2.65)$ & $(3.22)$\\
$k_F$ & $0.75$ & $1.12$ & $1.59$ & $1.89$ & $2.19$ & $2.47$ & $2.57$\\
\bottomrule
\end{tabular}
\end{table}
% 上句未完，接原书 p.77（PDF p092，不属本文件范围）。
